\section{Related Work}

Document parsing has rapidly evolved into a core research direction at the intersection of vision, language, and structure understanding, with the goal of converting heterogeneous digital and scanned pages into faithful, machine-readable representations. Based on how the parsing capability is trained and assembled, prevailing approaches can be broadly grouped into three families: pipeline-based methods, end-to-end methods, and reinforcement learning (RL)-based methods.

\subsection{Pipeline-based Methods}

A long-standing line of work follows a pipeline-based paradigm. Systems such as MinerU2.5~\cite{niu2025mineru2} and PaddleOCR-VL~\cite{cui2025paddleocr} decompose document parsing into a sequence of specialized stages: a layout analyzer first detects and classifies semantic regions (text blocks, tables, formulas, and figures), after which task-specific expert models recognize the content of each region in parallel, and an optional reading-order module reassembles the page into a structured format such as Markdown or HTML. This modular design offers favorable engineering ergonomics, supports targeted optimization of individual components, and enables high-throughput inference on standardized layouts. However, the very staging that grants modularity also introduces error accumulation along the pipeline: imperfect region proposals or misclassified blocks propagate to downstream recognizers and reading-order predictors, where they are difficult to recover. The reliance on hand-engineered interfaces between stages further limits adaptability to documents whose layouts deviate from the assumptions baked into each module.

\subsection{End-to-end Methods}

End-to-end approaches instead seek to subsume the full parsing stack within a single model. Representative systems such as the dots.ocr series~\cite{dotsocr2025,zheng2026multimodalocrparsedocuments} and the DeepSeek-OCR series~\cite{wei2025deepseek,wei2026deepseek2} fine-tune large Vision-Language Models (VLMs) via supervised fine-tuning (SFT) to map a document image directly to its structured representation, jointly handling text recognition, layout grounding, table and formula structuring, and reading-order reconstruction. By learning holistic visual--textual representations, these models avoid hand-crafted inter-stage interfaces and achieve strong in-distribution accuracy. Their effectiveness, however, is tightly coupled to the diversity and fidelity of the training corpus: the scarcity of large-scale, faithfully annotated parsing data (covering heterogeneous layouts, fine-grained element semantics, and globally consistent reading order) remains a fundamental bottleneck, and SFT-only models often degrade noticeably on out-of-domain (OOD) templates, low-resource languages, or specialized verticals such as charts and chemical formulas.

\subsection{RL-based Methods}

To move beyond the limits of token-level imitation, a growing body of work explores reinforcement learning as a post-training strategy for document parsing. Approaches such as~\cite{wang2025infinity, chen2025logics, poznanski2025olmocr2, ling2025table2latex} optimize VLMs against outcome-oriented rewards that reflect the quality of the parsed output, and have demonstrated encouraging gains across document types. Nevertheless, most existing RL pipelines remain narrowly scoped: they optimize a single, predominantly textual reward signal and treat related sub-capabilities (element parsing, table and chart decoding, chemical formula recognition, and document-level question answering) as isolated heads or sequential stages, leaving structural fidelity, spatial alignment, and cross-task transfer under-exploited. In contrast, our Infinity-Parser2 introduces a verifiable, multi-task reward system that drives Joint Reinforcement Learning across eight co-trained objectives (document parsing, layout analysis, table parsing, math formula parsing, chart parsing, chemical formula parsing, document VQA, and general multimodal understanding), unifying perception, structure, and reasoning within a single end-to-end optimization signal and thereby promoting robust generalization across heterogeneous documents and downstream tasks.